\subsection{Basic Properties of Logarithms} 
Even though the logarithm is only defined as an inverse, we figure out some basic properties.
\begin{description}
\item[Domain] The domain of $f(x)=\log_b(x)$ is the same as the \makeblank[1.25in] of\linebreak {$f\inv (x)=b^x$}, \makeblank[1in].
\item[Range] The range of $f(x)=\log_b(x)$ is the same as the \makeblank[1.25in] of {$f\inv (x)=b^x$}, \makeblank[1in].
\end{description}
More generally, the domain is wherever the \makeblank[1.5in] of the log is \makeblank[1in], and the range is always \makeblank.
\begin{Example}
Find the domain and range of $f(x)=\log_3(2x-5)+4$
\begin{lpic}{}{5}
\end{lpic}
\end{Example}
We can also give some basic values of logarithms. All of these can be found by rewriting in exponential form, but they may be worth memorizing. These are true for any number $b$ where the log is defined: \makeblank.
\begin{lpic}{}{5}
	\twocol{-1}{5}{Exponential Fact}{Logarithm Fact}
	\twocollblleft{-2}{$b^0=1$}{$\log_b(1)={}$}
	\twocollblleft{-3}{$b^1=b$}{$\log_b(b)={}$}
	\twocollblleft{-4}{$b^{-1}=\tfrac{1}{b}$}{$\log_b\left(\tfrac{1}{b}\right)={}$}
	\twocollblleft{-5}{$b^n=b^n$}{$\log_b{b^n}={}$}
%	\twocollblleft{-6}{$b^n=\left(\frac{1}{b}\right)^{-n}$}{$\log_{1/b}(x)={}$}
\end{lpic}
\begin{Example}
Find each log.
\begin{itemize}
\item $\log_4(1)=\makeblanksmall$
\item $\log_4(4)=\makeblanksmall$
\item $\log_4\left(\tfrac{1}{4}\right)=\makeblanksmall$
\item $\log_4\left(4^{3/2}\right)=\makeblanksmall$
\end{itemize}
\end{Example}